\subsection{THETA3-A13: the actual finite arrays, their spans, and tame globalization}
\label{A13}

This section gives restrictions and positive constructions for the
\emph{prescribed} column arrays.  It does not replace a prescribed
\texttt{xiAt} by an arbitrary Schwartz function.  The finite-field and
idele arguments below are new proofs.  Literal archive definitions and
the named archive existence theorem are source statements; no Lean
build or transitive dependency certification is asserted.

\paragraph{The defining objects and the coefficient field.}
All paths in this paragraph are relative to the pinned source root.
The literal global definition is
\texttt{Development/AbelianVarietyTheory/Definitions/AutomorphicMembership/}
\texttt{XiAtPenultimateEqXiCpl1.lean}, lines 146--162.
It uses \texttt{XiEntry} when $i+2<m$, \texttt{XiCpl1}
when $i+2=m$, and \texttt{xiLast} when $i+1=m$.
The last is the displayed \texttt{xiOver}; all three carry
the same \texttt{phiT} factor.
The notation \texttt{TC} is an abbreviation for \texttt{ColIn},
in \texttt{LatAtLattice.lean}, line 40, in that same directory.
The latter structure is in
\texttt{TrueColumnChoicesHeckeCharReferenceEqOne.lean},
lines 165--245: it has a finite injective family of tuples,
the reference tuple $1$, the specified reference characters,
local constants, local double-coset indices and their tying laws.
Thus a \texttt{TC} is more than a list of local characters.

The local definitions are in
\texttt{Formalism/LocalToGlobal/Columns/FiniteSumNonzeroLargeDepth.lean}:
\texttt{phiT}, lines 646--648;
\texttt{PhiQ}, lines 711--714;
\texttt{PhiG}/\texttt{PhiGS}, lines 780--805;
\texttt{om}, lines 553--556;
\texttt{xx1}/\texttt{E1}, lines 568--579;
\texttt{PhiAt}/\texttt{XiAt}, lines 927--948;
\texttt{XiEntry}/\texttt{XiCpl1}, lines 1116--1138.
The last local factor is \texttt{locTop}, lines 391--393 of
\texttt{Development/AbelianVarietyTheory/Definitions/GlobalTwisting/}
\texttt{IntegralGLSuccEmbeddingMemParahoricTimesPrincipa.lean}.

Write $K$ for the coefficient field, to distinguish it from the
imaginary quadratic subfield of the CM host.  The literal
\texttt{Fields.K} is a number field
(\texttt{AbelianHostFieldConditions.lean}, lines 19--44).
For the actual base it is $B.\mathrm{LL}$
(\texttt{BaseDataEmbeddingConj.lean}, lines 13--20).
There is an additional, concrete restriction on a fixed $K$.
In the actual canonical case \texttt{IsTw}, the proof
\texttt{root\_of\_ok} in \texttt{IsTwIffCanonEq.lean}
shows that every new twisting prime of residue cardinality $q$
requires a primitive $q^{2m}$th root of unity in $K$;
\texttt{sqrt\_of\_ok} requires a square root of $q$.
The proof uses the stipulated local additive character at depth
$2m$, its tying to the global character, and the fact that the
reference $\xi_0$ is a unit at this prime.  A fixed number field
therefore admits only finitely many such rational primes.
A12's choice $t=1$ removes a large power from its isolating
congruence, but not this separate canonical-depth requirement.

\paragraph{Two unconditional restrictions on the actual arrays.}
For a valid field and prime datum, put $f=[F:\Q]$.
The source has
\[
        |v_T(u)|=f>0,\qquad |v_G(j)|=f>0.
        \tag{A13a}
\]
These are \texttt{Primes.Ok.vT\_card} and
\texttt{vG\_card}, lines 216 and 231 of
\texttt{AbelianHostFieldConditions.lean}; positivity follows
from \texttt{Fields.Ok.card\_Tau} and $[F:\Q]>0$.
Every valid decoration has a nonzero residue functional at
each $T$-prime, by \texttt{Deco.Ok.f\_ne}, line 88 of
\texttt{StdLatUnimodular.lean}.
Consequently
\[
                         \mathrm{xiAt}(0)=0.
                         \tag{A13b}
\]
Indeed, \texttt{phiT} requires $f'_P\circ x=f_P$ on the
localized source lattice.  At $x=0$ its left side is zero:
$f'_P(0)=0$ follows from the displayed scalar-linearity
condition with scalar zero.  Its right side is nonzero on
some lattice vector.  Thus $\mathrm{phiT}(0)=0$ at any one
$T$-prime.  All branches of the global array contain this
factor, or are identically zero.

There is also the stronger rank restriction
\[
 \mathrm{xiAt}_i(x)\ne0
           \quad\Longrightarrow\quad
             \operatorname{rank}_E(x)=N(i),
       \qquad N(i)=k+i(t+1),\quad 0\leq i<m.
 \tag{A13c}
\]
Here the target dimension is $N(i+1)$.
To prove this directly, choose any one $G$-prime using
(A13a).  Set $c_i=m-1-i$ and $r=k+m-1$.
In \texttt{PhiG} and \texttt{PhiGS} at step $i$,
the nonzero reciprocal determinant factor is the determinant
of a square block of size
\[
               r+it=N(i)+c_i.
\]
Its reduction is a unit whenever the local function is nonzero.
The coordinate operation \texttt{entB} is an invertible row
and column permutation followed by nonzero diagonal scalings;
the displayed piecewise insertion \texttt{fIdx} has inverse
\texttt{fPos}.  Hence the underlying enlarged matrix has full
column rank $N(i)+c_i$.  The matrices in \texttt{om} are
invertible and do not change this conclusion.

For $i+2<m$, each summand of \texttt{E1} has matrix
\texttt{xx1}(x,t).  Its first $N(i)$ columns are the columns
of $x$ with zero auxiliary rows, and only $c_i$ columns
remain.  Thus
\[
       \operatorname{rank}(\mathrm{xx1}(x,t))
                     \leq\operatorname{rank}(x)+c_i.
\]
If $x$ is not injective, every summand is zero, so is
\texttt{XiAt}, and so is the global product.
For $i+2=m$, one has $c_i=1$ and
\texttt{joinCol}(x,y) has rank at most
$\operatorname{rank}(x)+1$, giving the same conclusion.
For $i+1=m$, one has $c_i=0$; the last local factor is
\texttt{PhiTop} applied to a row projection of an invertible
coordinate change of $x$.  Its required unit minor has size
$N(i)$, and a row projection cannot increase rank.
This proves (A13c) in all three cases.  In particular, the
actual array span is a proper subspace of all finite arrays;
this statement does not imply that its first-derivative image
is a proper subspace.

\paragraph{The exact incidence span at a $T$-prime.}
Let $\kappa$ have $q$ elements, let
$M=\operatorname{Mat}_{b\times a}(\kappa)$, and fix a nontrivial
additive character $\psi$ of $\kappa$.
For nonzero row vectors $v\in(\kappa^b)^*$ and
$u\in(\kappa^a)^*$, put
\[
                    I_{v,u}(X)=1_{\{vX=u\}}.
                    \tag{A13d}
\]
Their complex span is exactly
\[
 \{F:M\to\C:\ F(0)=0,
       \ \widehat F(A)=0\text{ whenever }\operatorname{rank}(A)>1\}.
 \tag{A13e}
\]
The Fourier pairing is $\psi(\operatorname{Tr}(A^tX))$;
using its inverse only reverses signs and leaves the assertion
unchanged.  Its dimension is
\[
                 \frac{(q^a-1)(q^b-1)}{q-1}.
                 \tag{A13f}
\]

Here is a full proof.  For a fixed projective row $[v]$,
the functions $F(vX)$ with $F(0)=0$ form a space of
dimension $q^a-1$, spanned by its nonzero-fibre indicators.
Their Fourier support is contained in the frequency block
$\{v^tt:t\in\kappa^a\}$: summing an additive character over
the kernel of $X\mapsto vX$ gives zero unless the frequency
annihilates that kernel.  Distinct projective rows have disjoint
nonzero frequency blocks, and these blocks partition all
nonzero rank-one matrices.  The sum of these spaces is direct.
Indeed, in a zero sum each summand has all nonzero Fourier
coefficients zero, so is constant; its value at zero makes
that constant zero.  There are $(q^b-1)/(q-1)$ projective
rows.  The entire Fourier-rank-at-most-one space has one more
dimension, its zero frequency included; evaluation at zero
is a nonzero functional on it.  The direct sum just constructed
has the dimension of its kernel, proving (A13e)--(A13f).

This is a literal local description of \texttt{phiT}, not
merely an analogy.  Keep both lattices standard and their
$G$-frames and $Q$-frames fixed.  Reduction at $P$ identifies
the source and target residue modules with $\kappa^a$ and
$\kappa^b$.  Every chosen nonzero $u,v$ gives a valid
decoration: define its functional by that residue-linear
map on the localized lattice and by zero off it.  Additivity,
scalar-linearity and the off-lattice condition are immediate;
nonvanishing on the global standard lattice follows from
surjectivity of $\OO_E\to\kappa$.  The same construction
works independently at every $T$-prime.  The matrix-integrality
part of \texttt{phiT} is kept as well.  Thus, on integral
residue matrices, the tensor product of the spaces (A13e)
is precisely the span of the local \texttt{phiT} factors
obtained by varying these raw decorations.  The corresponding
global-array span is the image of this space under multiplication
by the other fixed local factors, followed by the stated class
equivariance.  Equality with an unmultiplied space is asserted
only on a CRT slice where those other factors have fixed nonzero
values.
Passing to the chosen class representatives still uses the
archive's equivariance maps; no new list of representatives
or new full \texttt{TC} is being silently declared.

These restricted incidence functions can isolate a finite
window even though they cannot isolate every residue matrix.
Let $\mathcal W\subset M$ be finite, with $X_0\in\mathcal W$
nonzero.  If $q>|\mathcal W|$, there is $v\ne0$ such that
\[
 vX_0\ne0,\qquad
 I_{v,vX_0}(X)=0\quad(X\in\mathcal W,\ X\ne X_0).
 \tag{A13g}
\]
Each forbidden equation $vX_0=0$ or $v(X-X_0)=0$
is a proper linear subspace of the row space, of cardinality
at most $q^{b-1}$.  At most $|\mathcal W|$ such subspaces
are forbidden.  Their union has cardinality less than $q^b$.
Choosing a row outside it proves the assertion and also
ensures that the source decoration $u=vX_0$ is nonzero.
For a finite set of rational lattice matrices, the same test
applies whenever their reductions are defined, distinct from
$X_0$, and $X_0$ has nonzero reduction.  These are exact
finite divisibility tests.  If several $T$-primes are present,
one can impose (A13g) at one of them and merely preserve
$X_0$ at the others, provided its reduction there is nonzero.

\paragraph{What tame determinant characters and matrix translates span.}
There is an equally explicit answer for the first-step
\texttt{PhiQ} seed.  Before its Jacquet-type finite summation,
its source dimension is $r=k+m-1$.  After the displayed row
permutation, its reduction on
$\operatorname{Mat}_{b\times r}(\kappa)$ is
\[
 F_\chi(X)=
 \begin{cases}
   \chi(\det X_0),&X=\binom{X_0}{0},\quad X_0\in\mathrm{GL}_r(\kappa),\\
   0,&\text{otherwise}.
 \end{cases}
 \tag{A13h}
\]
Here $\chi$ is a multiplicative character of $\kappa^\times$.
Let $p(X)$ be the wedge of the $r$ columns.  The span of
all left and right general-linear translates of (A13h),
as $\chi$ ranges over all multiplicative characters, is
exactly the pullback under $p$ of all functions on the
nonzero decomposable $r$-vectors; these functions are extended
by zero on matrices of rank less than $r$.

Indeed $\mathrm{SL}_r(\kappa)$ acts transitively on the
bases of an $r$-plane with any specified exterior product.
Every function (A13h), and every translate of it, is invariant
under this right action.  Conversely, a left translate can
place the image plane at any specified $r$-plane $W$.
Choose a nonzero $w\in\bigwedge^rW$.  On that plane its
functions are exactly $\chi(d)$ for $p(X)=dw$.
Character orthogonality gives
\[
 1_{\{d=d_0\}}
       =\frac1{q-1}\sum_{\chi}\chi(dd_0^{-1}).
 \tag{A13i}
\]
This proves the span claim.  If the trivial character is
excluded, the exact span consists of the functions whose
sum over the $q-1$ nonzero multiples of each $w$ is zero.
In particular
\[
 1_{\{d=d_0\}}-1_{\{d=d_1\}}
   =\frac1{q-1}\sum_{\chi\ne1}
           (\chi(d_0)^{-1}-\chi(d_1)^{-1})\chi(d)
 \tag{A13j}
\]
has coefficient absolute sum at most $2$.
Thus a second residue can cancel the trivial component
exactly, but the result is a delta on an \emph{oriented
image plane}, not on an individual full-rank matrix.

The source's first local \texttt{XiAt} retains a corresponding
restriction even after the finite summation.  For a rational integral matrix
$A\in\mathrm{GL}_k(E)\cap\mathrm{GL}_k(\OO_{E,P})$, its defining formulas give
\[
       \Xi_{0,P}(xA)=\chi(\operatorname{red}_P(\det A))\Xi_{0,P}(x).
       \tag{A13k}
\]
For clarity this identity is in the displayed local coordinates;
the locally constant identity extends to the completed local
matrix group, and integral changes of source frame preserve its
$\mathrm{SL}_k(\OO_{E,P})$-invariance.
At step zero, \texttt{ChIdx.Ok.zero\_h} sets the source
translate to $1$.  The enlarged source has precisely $r$
columns.  Right multiplication by
$\operatorname{diag}(A,1_{m-1})$ preserves the integrality
and the zero lower block of (A13h) and multiplies its unit
minor by $\det A$.  Also
\[
 \mathrm{xx1}(xA,t)=\mathrm{xx1}(x,t)
                    \operatorname{diag}(A,1_{m-1}),
\]
while its auxiliary character \texttt{chArg} is unchanged.
Every summand and its normalization therefore obey (A13k).
When $m=2$, the first step is \texttt{PsiAt} instead;
the identical calculation uses
$\mathrm{joinCol}(xA,y)=\mathrm{joinCol}(x,y)
\operatorname{diag}(A,1)$.
This proves (A13k), including right special-linear invariance,
without assuming that arbitrary local translates are admitted
by the actual finite index set.  The exact seed span (A13i)
does not by itself prove surjectivity after those index and
Jacquet constraints.

\paragraph{Exact tame anticyclotomic globalization.}
Assume $E/F$ is CM quadratic and unramified at every finite
place.  Let $\mathcal S$ be a finite set of split finite
places of $F$, with a selected prime $P_v$ of $E$ above each
$v$.  Write $\bar P_v$ for its conjugate and $\kappa_v$
for the common residue field.  Let
\[
 R_{\mathcal S}=\prod_{v\in\mathcal S}\kappa_v^\times,
 \quad
 U(\mathfrak m)=
 \prod_{P\nmid\mathfrak m}\OO_{E,P}^\times
 \prod_{P\mid\mathfrak m}(1+P\OO_{E,P}),
 \quad
 \mathfrak m=\prod_{v\in\mathcal S}P_v\bar P_v.
\]
The finite group relevant here is
\[
 C_{\mathcal S}=
 \A_{E,f}^\times/(E^\times\A_{F,f}^\times U(\mathfrak m)).
 \tag{A13l}
\]
It is a quotient of the usual finite ray class group.
For this standard background see Milne,
\emph{Class Field Theory}, v4.03 (2020), Theorem V.1.7,
printed pp.~150--151, and Proposition V.4.6,
printed p.~173; both statements and their complete displayed
proofs were read for this comparison.

At a split place use the ratio convention
$(a,b)\mapsto\bar a\bar b^{-1}\in\kappa_v^\times$,
where bars on $a,b$ in this sentence mean residue reduction.
Local unit representatives give a well-defined homomorphism
$R_{\mathcal S}\to C_{\mathcal S}$.  Its kernel is exactly
\[
 \Delta_{\mathcal S}\mu(E)
    =\{(\zeta\bmod P_v)_v:\zeta\in\mu(E)\}.
 \tag{A13m}
\]
Here is a proof that includes the possible global-unit issue.
If a local-unit idele represents zero in (A13l), write it
as $a f w$ with $a\in E^\times$, $f\in\A_{F,f}^\times$,
$w\in U(\mathfrak m)$.  The valuations of $a$ at conjugate
primes are equal, so $a/\bar a$ is a global unit.  Every
complex conjugate of this unit has absolute value one.
It is therefore a root of unity: the powers are algebraic
integers of bounded degree whose conjugates all have
absolute value one; their monic minimal-polynomial
coefficients belong to fixed finite sets of integers,
so the set of powers is finite.  Taking the local ratios
in the chosen order gives precisely its image in (A13m).

Conversely, for every $\zeta\in\mu(E)$ one can choose
$a\in E^\times$ with $a/\bar a=\zeta$.
For $\zeta\ne-1$ take $a=1+\zeta$; for $\zeta=-1$
take any nonzero purely imaginary element of $E$.
Its valuations at conjugate primes agree.  Since every
finite ramification index in $E/F$ is one, choose an
$F$-idele $f$ with valuation $-v_P(a)$ at the underlying
place.  Then $af$ is an $E$-unit idele at every finite
place.  Let $u$ equal $af$ at the places in $\mathcal S$
and $1$ elsewhere.  The quotient $u/(af)$ lies in
$U(\mathfrak m)$, and the ratio of the two components
of $u$ at $v$ reduces to $\zeta\bmod P_v$.
Thus its class is zero, proving (A13m).  The ramification
index condition is used exactly in the choice of $f$;
no principal-ideal or class-number-one assumption was used.

It follows that prescribed characters
$\lambda_v:\kappa_v^\times\to\C^\times$ occur as the tame
inertia restrictions of a finite-order character trivial
on $E^\times\A_{F,f}^\times$ if and only if
\[
              \prod_{v\in\mathcal S}
                   \lambda_v(\zeta\bmod P_v)=1
                         \quad(\zeta\in\mu(E)).
 \tag{A13n}
\]
Necessity follows from (A13m).  For sufficiency, the
product character descends to the image of $R_{\mathcal S}$
in $C_{\mathcal S}$.  Every character of a subgroup of
a finite abelian group extends to the group.  One elementary
proof adjoins one group element at a time: if $x$ has order
$d$ modulo the old subgroup, choose a $d$th root of the
specified value on $x^d$, and set the new character on
$hx^j$ equal to the old value on $h$ times that root to
the power $j$.  This is well-defined and enlarges the
domain.  Iteration terminates.  The final character has
values in roots of unity of order dividing
$|C_{\mathcal S}|$.  Its pullback to finite ideles has
the requested restrictions, is trivial on all wild inertia,
and is unramified outside $\mathcal S$.
The class group influences the choices and possible orders
of the extension, not an additional condition on (A13n).

This gives an actual \texttt{Tup} once all places in
$\mathcal S$ are split degree-one places outside the
forbidden set: compare \texttt{GoodPl}, \texttt{finChar}
and \texttt{tupGrp}, lines 64--115 of
\texttt{Formalism/LocalToGlobal/Divisibility/Tuples.lean}.
The exact residue character at $(a,b)$ is
$\lambda_v(\bar a\bar b^{-1})$; changing the selected
prime inverts this convention.

\paragraph{One auxiliary prime removes the root-of-unity condition.}
Prescribe any finite collection of local characters at the
main places.  Choose one additional allowed split degree-one
place $v_*$ whose residue characteristic does not divide
$|\mu(E)|$.  Reduction injects $\mu(E)$ into
$\kappa_{v_*}^\times$: roots of unity of order prime to
the residue characteristic cannot lie nontrivially in
$1+P_{v_*}$, as follows either by the binomial expansion
of $\zeta^n=1$ with $n$ a unit or by uniqueness of simple
Hensel lifts.  The inverse of the main product in (A13n)
is a character of that embedded subgroup.  Extend it to
$\kappa_{v_*}^\times$ by the same finite-group argument.
Then (A13n) holds and the prescribed main characters
globalize without alteration.  One auxiliary prime works
for the entire finite collection, with a possibly different
auxiliary character for each member.

For a finite forbidden set such primes are available by
Chebotarev, choosing rational primes split completely in
a Galois closure of $E$ and excluding the finite list.
The exact background locator is Milne, \emph{Algebraic Number
Theory}, v3.08 (2020), Theorem 8.31 and Corollary 8.32,
printed p.~147; these give the positive density of the
identity Frobenius class and hence infinitude.
For a \emph{fixed} actual \texttt{TC}, however, its set
$I.S$ is literally an arbitrary set of rational integers;
\texttt{ColIn.Ok} only requires it to contain the specified
$T$-primes, the two $G$-primes, and the primes dividing
the host conductor.  Thus the assertion for that fixed
datum requires an exhibited spare allowed place; it does
not follow solely from the type of \texttt{TC}.

The auxiliary prime changes its local \texttt{PhiQ}/\texttt{XiAt}
factor.  Its value and support are part of the new input.
In particular, the exact character identity (A13j)
at a main prime is not a global-array identity until
these auxiliary factors are included.  There are only
$|\mu(E)|$ possible restrictions to cancel, so one can
choose at most that many auxiliary residue characters
for this purpose.  This is a finite, explicit local cost,
not a proof that the auxiliary kernel equals $1$.

\paragraph{Coefficient enlargement and the remaining finite comparison.}
For the finite character family, a single number field
$K'$ can contain its values, all required
$\zeta_{q^{2m}}$, and all required $\sqrt q$.
Take the compositum of the old normal coefficient field
with the corresponding finite cyclotomic and quadratic
extensions, in the chosen complex embedding.  Complex
conjugation restricts to the old involution.  The
\texttt{Fields.Ok} conditions transport: the $E/F$ data
are unchanged; every embedding of $E$ into the new field
already has image in the old normal field, so the list
of embeddings and the defining equalizer $K_q\subset E$
are unchanged; total complexity and normality persist.
This is a statement about the coefficient package, not
an alteration of the rigid full CM host field or its
Galois frame.

In fact the entire literal \texttt{Base} survives this
coefficient enlargement.  Its complete definition is
\texttt{Development/AbelianVarietyTheory/Definitions/Foundations/}
\texttt{RootOfUnityEightQMemCoeffGaloisSubfield.lean},
lines 261--304.  Keep \texttt{Pre} unchanged.  Embed the
old $\mathrm{LL}$ in the chosen normal field $K'$ and
use the images of the \emph{old} intermediate fields
$\mathrm{LL}_0$ and $K_s$ as the new intermediate fields.
They are isomorphic to their originals, rather than
being replaced by their composita with $K'$.
The old $\mathrm{LL}_0$ is still Galois over $\Q$;
the inclusions $E\subset K_s\subset\mathrm{LL}_0$
and every stipulated root of unity transport unchanged.
An isomorphism of number fields identifies their rings
of integers, ideals and principal ideals.  It therefore
transports each generator in \texttt{capit} and
\texttt{kPsi}, with exactly the same principal-ideal
equalities.  Every complex embedding of the image of
$K_s$ is uniquely the old embedding composed with this
isomorphism's inverse, so the norm equalities quantified
over \emph{all} embeddings in \texttt{kPsi} also persist.
The roots required merely in $\mathrm{LL}$ remain in
$K'$.  Conjugation on the chosen normal subfield of
$\C$ gives the required involution, compatible with
the old one and the prescribed complex embedding.
These are all the additional clauses of \texttt{Base}.
In particular the CM host, its Galois frame, the ranks,
the $T$-primes and their fixed windows, and the old
integral coefficient ring $\OO_{K_s}$ are unchanged.

The canonical local indices are supplied by
\texttt{twIdx\_ok\_closed}, lines 87--102 of
\texttt{Formalism/LocalToGlobal/Columns/}
\texttt{UnfoldingNonvanishingDoubleCosetHolds.lean}.
Its hypotheses explicitly include the existing field,
prime, reference-character and Weil tying conditions,
the new good-prime and unit conditions, and the roots
and square roots just listed.  Its displayed proof was
read; it invokes the archive's nonvanishing and finite-index
theorems.  We do not replace those invocations by a claim
of a fresh independent proof.

\paragraph{Complete realization of a prescribed finite character family.}
Here the source provides a stronger positive conclusion
than merely the existence of the local indices.  Fix the
original \texttt{Pre}, an oriented \texttt{PreX}, the
reference frame $f$, a prime-selection map \texttt{sel}
and a valid ambient theta datum $J$.  Retain the displayed
diagonal-space and $\mathrm{duStar}$ identities of that
geometric input.  Let $S$ contain its $T$-primes,
the two $G$-primes and all primes dividing the host
conductor.  Suppose that the finite family $\Xi$ contains
$1,\delta,\delta'$ and that every $\xi\in\Xi$ is a
finite-order character trivial on $\A_F^\times$, tame
at every finite place, and ramified only at split,
absolutely unramified degree-one places outside $S$.
These are precisely the clauses of \texttt{Chr.IsTwist}
in
\texttt{Development/AbelianVarietyTheory/Definitions/AutomorphicMembership/}
\texttt{IsRationalUnitaryDeterminantMapRationalUnitaryTo.lean},
lines 111--114.  Thus the characters constructed above
satisfy them whenever the main and auxiliary places
have been selected as stipulated.

Assume the original coefficient field contains the
normalized reference-character values required by
$\mu_{0,j}.\mathrm{ValuesIn}(0)$ and
$\eta.\mathrm{ValuesIn}(1/2)$, and the square roots
of the two fixed $G$-primes.  These are explicit reference
conditions, satisfied by the original fully tied
column construction.  Enlarge only $\mathrm{LL}$ as
above to contain the values of all $\xi\in\Xi$ and
$\sqrt q,\zeta_{q^{2m}}$ for the finite set of their
ramified residue cardinalities.  This is finite: if
the character orders divide $M$ and $\mathcal Q$ is
that set of residue cardinalities, one possible
compositum has the coarse relative degree bound
\[
 [K':K]\ \leq\ \varphi(M)
                   \prod_{q\in\mathcal Q}2\varphi(q^{2m}).
 \tag{A13o}
\]
All factors on the right bound the degrees of the
explicit cyclotomic or quadratic extensions being
adjoined.  Intersections can only decrease this bound.
For the characters extended from a single
$C_{\mathcal S}$ above one may use $M=|C_{\mathcal S}|$,
also including the orders of the fixed reference
characters if necessary.

There then exists an actual finite \texttt{TC} on this
same geometric \texttt{Pre}, with exactly the character
set $\Xi$, the prescribed \texttt{sel} and the same
ambient datum $J$, satisfying \texttt{Ok},
\texttt{Laws} and \texttt{Law0}.  It may be chosen
canonical in the sense \texttt{IsTw}.  If the original
\texttt{Pre} has \texttt{StagesInhab.HasLetters}, its
\texttt{Later} data exist as well.

\textbf{Proof and exact source scope.}
The full constructor is \texttt{InEx.exists\_colIn},
lines 195--289 of
\texttt{Development/AbelianVarietyTheory/Proofs/IwasawaDecomposition/}
\texttt{ExistsColumnInputAssembled.lean}.
Its condition \texttt{Holds}, lines 126--128, is exactly
the value-field, square-root and depth-$2m$ root-of-unity
condition just arranged.  Its proof first constructs
the reference finite characters, reads each member of
$\Xi$ as a tuple, applies the local-index existence
theorem to each member, and uses $\Xi$ itself as its
finite injective index set.  It returns \texttt{Ok},
the asserted equalities of character range, $S$,
\texttt{sel} and $J$, and the exact constants
$M=A=B=m$, $D_Q=2m$.
In particular its \texttt{Deep} condition holds.
The complete proof of \texttt{canon\_ok} and
\texttt{canon\_isTw} in \texttt{IsTwIffCanonEq.lean}
then replaces only its constants and indices by their
canonical versions.  The ambient theta datum,
characters and selected primes are preserved.

For this actual valid input, the statements needed next
are proved theorems, not additional Prop declarations.
In
\texttt{Development/AbelianVarietyTheory/Definitions/AutomorphicMembership/}
\texttt{ExistsOkPins.lean},
\texttt{TrueColsTest.laws}, lines 86--103, proves
$I.\mathrm{Ok}\Rightarrow I.\mathrm{Laws}$;
\texttt{law0}, lines 152--159, proves
$I.\mathrm{Ok}\Rightarrow I.\mathrm{Law0}(l)$ for
every $l:I.\mathrm{Laws}$; and \texttt{later},
lines 194--204, constructs all five later stages
from $l$ and \texttt{HasLetters}.
The first assembles the chain, level, period and lattice
theorems.  The second pulls functions back through the
step-zero class bijection, whose representatives have
the same underlying lattice, and hence preserves the
integral lattices as well as their complex embeddings.
The third constructs the finite-dimensional row model
and its divisor, level, local and jet stages.
All three complete displayed proofs were read.
\texttt{HasLetters} depends only on the unchanged
\texttt{Pre}; in the archive's contextual construction
\texttt{JetEx.hasLetters\_of\_over}, lines 16--26 of
\texttt{ExistsColInJetsIn.lean}, supplies it from
the existing \texttt{Ctx.Over} hypothesis.
This proves the composite realization statement using
the named source theorems.\hfill$\square$

The stronger specialized \texttt{ColW} carries another
restriction.  Its literal definition, lines 23--36 of
\texttt{PositionOfIntro.lean}, fixes its character range
to $T.\mathrm{set}(f)$ for one \texttt{TwDat}.
The latter definition, lines 79--97 of
\texttt{CmQuadraticFormEqAlternatingForm.lean}, requires
the auxiliary character indexed by a coefficient prime
$p$ to have $p$-power order, with its \texttt{IsPsi},
conjugation and avoidance conditions.  Its range consists
of $1,\delta,\delta'$ and the specified products of
one of $\delta,\delta'$ with these auxiliary characters.
An arbitrary globally admissible finite family is not
automatically of this special form.
For this specialized witness the forbidden-prime set is,
however, finite.  The host's \texttt{levelPrimes} is
literally a finite set (\texttt{Formalism/LocalRingCardinality/}
\texttt{DataHost.lean}, line 69), and the complete proof
\texttt{finite\_offSet}, lines 208--216 of
\texttt{CmQuadraticFormEqAlternatingForm.lean}, places
all excluded primes in a finite divisor set together
with those level primes and one distinguished prime.
The \texttt{ColW.S\_eq} clause therefore proves that
its $S$ is finite.  Hence an auxiliary completely split
prime outside it, the main support, and the divisors of
$|\mu(E)|$ is available by the cited Chebotarev theorem.
The unresolved specialized-family condition is its
\texttt{TwDat} shape, not the existence of a spare prime.

There is also a separate positive arithmetic-jet source
theorem: \texttt{exists\_bound\_jetsIn}, lines 34--40
of \texttt{ExistsColInJetsIn.lean}, gives a finite list
of algebraic coefficient generators after which a valid
$I'$ with \texttt{JetsIn} exists and has the same
reference frame, forbidden set and character range.
Its conclusion does not retain the ambient datum $J$.
The following refinement retains $J$ by allowing that
finite generator list to depend on it.

\paragraph{A fixed-ambient-datum arithmetic-jet refinement.}
Under the preceding realization hypotheses and
\texttt{HasLetters}, after a further finite coefficient
extension the complete canonical \texttt{TC} can be
chosen with the \emph{same} $J$ and with
\texttt{JetsIn}.  One may simultaneously require the
same field to contain all normalized jets of any
prescribed finite list of total degrees, including
degree one, on the admitted lattices and torsion
arguments of these finitely many steps.  This is
an arithmetic-field assertion; it asserts neither
nonzero jets nor surjectivity of their span.

Here is the full refinement argument.  The crucial
source theorem is \texttt{JetEx.finAlg\_step},
lines 193--226 of
\texttt{Development/AbelianVarietyTheory/Proofs/IwasawaDecomposition/}
\texttt{DegreeEqTotDeg.lean}.  It applies to \emph{any}
valid ambient datum $J$ in an $I.\mathrm{Ok}$ input,
and to every total degree $n$.  It states that the
corresponding normalized jet set is contained in a
number field.  Its proof invokes the normalized-vacuum
algebraicity and finite-jet-set theorems, with the
displayed self-duality, parity, bad-prime and ambient
normalization hypotheses.  This source result is used
with those hypotheses, not inferred from arbitrary
analytic theta normalization.

Let $\mathcal B$ be the finite set of primes of $E$
above the fixed $T$- and $G$-primes and the ramification
of the finitely many characters in $\Xi$.
For a word $u$, a step $i<m$, and subsets
$S_0,R_0\subset\mathcal B$, put
\[
 M(u,R_0)=T(u)
             \left(\prod_{Q\in R_0}\mathrm NQ\right)^{2(m+1)}.
 \tag{A13p}
\]
The exponent is precise: the canonical
\texttt{twCst} has $A=B=m+1$, $M=m$,
as displayed in \texttt{Formalism/LocalToGlobal/Columns/}
\texttt{ChainTw.lean}, lines 217--226.
For each such finite index and each selected degree,
form the source's normalized \texttt{jetSet} using
$J$, $S_0$, modulus (A13p), and the fixed ideal
character $\chi$.  These sets are independent of the
coefficient field: \texttt{hostAmb} and
\texttt{jetSet\_amb}, lines 43--66 of
\texttt{DegreeEqTotDeg.lean}, prove this exactly.

Keep the members of this finite list which are
contained in a number field, and take their union.
It is contained in the compositum of finitely many
number fields.  Adjoin a finite set of algebraic
generators for that compositum to the character
coefficient field, using the same \texttt{Base}
enlargement construction.  This finite-choice argument
is the explicit mathematical content of the source's
\texttt{FinAlg} union step; it does not require a
decision procedure for testing the members of the list.

Now construct and canonicalize the actual input in
this enlarged field, retaining $J$.  For every one of
its words and steps, its actual sets
$\mathrm{Sb}(a,u)$ and $\mathrm{rowSG}(a)$ are subsets
of $\mathcal B$.  The complete proofs of
\texttt{Sb\_subset} and \texttt{rowSG\_subset},
lines 126--156 of the same file, establish precisely
this containment from $I.\mathrm{Ok}$ and the character
range.  Its actual modulus is (A13p), by the definitions
of \texttt{cAof}, \texttt{cBof} and \texttt{twCst}.
Apply \texttt{finAlg\_step} to this actual valid input.
Its jet set is consequently one of the retained sets
in the finite union, and all its values lie in the
chosen field.  Taking, for each $u,i$, the total degree
of the required embedding multidegree
$\nu\mathrm{Deg}(u,i,\tau)$ proves
\texttt{JetsIn}; the identity between these total
degrees is proved in \texttt{degree\_eq\_totDeg},
lines 75--84.  Adding any finitely many other degrees
to the initial list proves the additional assertion.
The universal laws and later-stage constructors apply
again without a change of \texttt{Pre}.\hfill$\square$

This proof improves the wrapper's preservation statement,
but only for a datum already satisfying the literal
\texttt{AmbIn.Ok}.  The ambient structure and that validity definition fix its
discriminant-normalized vacuum and differential generators;
it is not a statement that an arbitrary 032 normalization
is such a datum.  The comparison with a particular 032
section must still retain the actual section, tangent
and period factors, as in A8 and A12.  Nor does (A13o)
bound the degree of this further arithmetic-jet field:
the latter is supplied by the stated finite algebraic
generation theorem and its finite compositum.

Finally, (A13g) is an actual \emph{array}-window selector.
The precise source passage from this array to a holomorphic
theta section is \texttt{reads\_datum\_eq\_thC}, in
\texttt{Development/AbelianVarietyTheory/Proofs/DecoratedLatticeProofs/}
\texttt{DatumReadingSplitPlacesEqClassicalTheta.lean}.
Its complete proof expresses the output as a finite
Fourier sum of Heisenberg translates of the chosen vacuum;
the half-pairing phases, semicharacters, and lattice
representatives all occur in that sum.  Thus a finite
window in the raw $E$-matrix array is not automatically
the same window in the rational Schr\"odinger Gaussian
series of A12.  For any given full datum, use that explicit
finite transform and then the derivative matrix: apply it
to the incidence basis (A13d), or to the oriented-plane
character combinations (A13i)--(A13j), and compute its rank
and a nonzero minor with all common section and tangent
periods removed as in A12.  This is a smaller, exactly
specified finite comparison.  The new results above close
the global tame-character compatibility condition and
give literal array selectors and invariant subspaces;
they do not assert the unproved surjectivity of this
final transformed derivative matrix.

In the requested Grothendieck framework this last distinction
is essential: the finite theta-module comparison belongs
to the period and correspondence comparison discussed in
\emph{cote44}, scans 7--8.  It is not an algebraic Lefschetz
inverse or a cohomological projector of the kind discussed
in \emph{cote16}, scans 43--46.  A later inverse matrix is
algebraic only after its stated period factors have been
removed, and an actual correspondence still requires the
cohomological identification.
